\begin{abstract}
Large language models increasingly rely on external tools to access up-to-date information, perform computation, and interact with the outside world. For autoregressive models, tool use naturally fits generation: emit a call, wait for the result, and continue. Diffusion language models (dLLMs), however, repeatedly revise many regions in parallel, making a stop--call--resume interface unnecessarily restrictive.

We introduce \emph{Continuous Interaction Diffusion} (CID), a diffusion-native model--runtime architecture that integrates tool interaction into iterative denoising. CID separates an externally controlled fact channel, a thought channel represented by a \emph{Typed Cognitive Tensor}, and a display channel. Information needs can form persistent bindings to external sources; returned observations are projected into evolving cognition and can revise earlier thought and display regions. Static values can be re-projected without repeated I/O, while changing sources can be refreshed. The runtime can continue useful source-independent refinement while external work is in flight and selectively reopen affected state when new evidence arrives.

We formalize CID and instantiate it in trained diffusion language models, including its representation, asynchronous runtime, and training objectives. We evaluate the resulting models on read-only tool-augmented reasoning tasks with delayed, changing, and reusable external information, measuring end-to-end task quality, tool-use reliability, asynchronous efficiency, and revision under new evidence.
\end{abstract}
